\chapter{Giới thiệu}

\section{Bối cảnh và động lực}
Log hệ thống là nguồn dữ liệu quan trọng để giám sát vận hành, phát hiện sự cố
và hỗ trợ AIOps. Các hệ thống phân tán và siêu máy tính có thể sinh hàng triệu
bản ghi log trong thời gian ngắn. Do tính bán cấu trúc, phi dừng, mất cân bằng
lớp và tốc độ sinh dữ liệu lớn, phát hiện bất thường dựa trên log đòi hỏi cả
pipeline tiền xử lý phù hợp lẫn giao thức đánh giá có thể tái lập.

\section{Vấn đề nghiên cứu}
Các phương pháp phát hiện bất thường dựa trên log đã phát triển từ học máy cổ
điển, học sâu tuần tự hoặc Transformer đến mô hình ngôn ngữ lớn kết hợp truy
xuất. Tuy nhiên, kết quả công bố thường phụ thuộc dataset, parser, cách chia
dữ liệu, ngữ nghĩa nhãn và phần cứng. Khoảng trống của nghiên cứu là thiếu một
benchmark thống nhất cho ba họ phương pháp đại diện.

\section{Mục tiêu và phạm vi}
Luận văn đánh giá ba phương pháp đại diện: LightAD cho học máy cổ điển,
LogFormer cho Transformer và LogRAG cho mô hình ngôn ngữ lớn kết hợp RAG.
So sánh được thực hiện trên HDFS\_v1 và BGL trong một giao thức thống nhất.
Phạm vi bao gồm hiệu quả phát hiện, chi phí tính toán, đánh đổi hiệu quả--chi
phí, mẫu FP/FN và khả năng giải thích.

\section{Câu hỏi nghiên cứu}
\begin{itemize}
  \item RQ1: Các phương pháp đạt Precision, Recall, F1, PR-AUC và FPR như thế
  nào trên các dataset đã định nghĩa?
  \item RQ2: Chi phí huấn luyện, suy luận, bộ nhớ và API là bao nhiêu?
  \item RQ3: Đánh đổi giữa hiệu quả phát hiện và chi phí tài nguyên như thế
  nào?
  \item RQ4: Những kiểu FP/FN nào lặp lại và phương pháp nào cung cấp lời giải
  thích hữu ích hơn?
\end{itemize}

\section{Đóng góp}
\begin{enumerate}
  \item Xây dựng pipeline thu thập, parsing, gom nhóm và đánh giá có thể tái lập.
  \item So sánh ba họ kỹ thuật trong cùng điều kiện dataset và phần cứng.
  \item Phân tích lỗi và khả năng giải thích ở cấp mẫu, template và nhóm.
  \item Đưa ra khuyến nghị lựa chọn phương pháp theo ràng buộc vận hành.
\end{enumerate}

\section{Cấu trúc luận văn}
Chương 2 trình bày nền tảng dữ liệu log và các họ phương pháp. Chương 3 mô tả
thiết kế nghiên cứu. Chương 4 báo cáo kết quả. Chương 5 thảo luận, hạn chế và
hướng phát triển. Chương 6 kết luận.
